\documentclass[10pt,a4paper]{amsart}
\usepackage[margin=19mm]{geometry}
\usepackage{amsmath,amssymb,mathtools}
\usepackage[scr=rsfs,cal=euler]{mathalfa}
\usepackage{xfrac}
\usepackage[unicode,hidelinks,bookmarks=false]{hyperref}
\newcommand{\ie}{i.e.\ }
\newcommand{\cf}{cf.\ }
\newcommand{\resp}{respectively\ }
\newcommand{\Subsection}{\subsection}
\providecommand{\nobreakdash}{\nobreak\hbox{-}}
\setlength{\emergencystretch}{2em}
\multlinegap0pt
\usepackage{bm}
\usepackage{amssymb}
\usepackage{xcolor}
\usepackage{algorithm}\usepackage{algpseudocode}
\newtheorem{lemma}{Lemma}
\newtheorem{proposition}{Proposition}
\title{\bf Convergence of the Immersed Interface Method in Linear Elasticity}
\author{Sabia Asghar, Qiyao Peng, Etelvina Javierre, Fred J. Vermolen}
\date{Source: arXiv:2410.10436v2 (CC BY 4.0)}
\begin{document}
\maketitle

\noindent\emph{Source and licence.} \bf Convergence of the Immersed Interface Method in Linear Elasticity. Version arXiv:2410.10436v2 (CC BY 4.0), distributed by arXiv under the Creative Commons Attribution 4.0 International licence, http://creativecommons.org/licenses/by/4.0/, as recorded in the arXiv licence metadata for this exact version and shown on the abstract page https://arxiv.org/abs/2410.10436v2. Full licence notice: \texttt{licenses/arxiv-2410.10436-CC-BY-4.0.txt}.\par\smallskip

\noindent\textit{Editorial scope.} Counted unit: the article's proposition corr2 (source0.tex lines 419--422). The article's complete original proof of it is delivered with it (source0.tex lines 423--443, one \texttt{proof} environment of 21 lines). No mathematics is written, repaired or re-derived: the statement and the proof are the article's own lines, in the article's own notation, and no retained line is substituted or re-flowed.  Retained uncounted as the source-local context the proof needs: lines 196--201; lines 203--207; lines 299--314.  Nothing was omitted from any retained span, so the package carries no omission record.  No retained line contains a citation, so the wrapper carries no bibliography and no bibliography entry.  No retained line contains a figure environment, an image-inclusion command or drawing code, so no image is delivered and the package declares no asset dependency.  Editorial formatting only, in addition: an AMS \texttt{amsart} preamble that restates the article's own declarations and macros used by the retained lines, namely the environments \texttt{lemma}, \texttt{proposition} on separate counters, together with the standard packages those lines need.  The retained environments print in this document as Lemma 1, Proposition 1 in order of appearance, whereas the article prints the same items with the numbering of its own full document.  The complete native source of the article as distributed in the arXiv e-print for this exact version is bundled unchanged as \texttt{sources/166796/source0.tex}.
\begin{equation}
\boldsymbol{u}^{\ast }(\boldsymbol{x})=\int_{\Gamma }Q(\boldsymbol{x}%
^{\prime })\boldsymbol{G}(\boldsymbol{x},\boldsymbol{x}^{\prime })%
\boldsymbol{n}(\boldsymbol{x}^{\prime })d\Gamma (\boldsymbol{x}^{\prime
}), \qquad {\boldsymbol x} \in \mathbb{R}^d,  \label{eq:u_star}
\end{equation}

\begin{equation}
{\boldsymbol{u}}_{h}^{\ast }({\boldsymbol{x}})=\sum\limits_{j=1}^{m} w_j Q(%
\boldsymbol{x}_{j})\boldsymbol{G}(\boldsymbol{x},\boldsymbol{x}_{j})%
\boldsymbol{n}(\boldsymbol{x}_{j}), \qquad {\boldsymbol x} \in \mathbb{R}^d.  \label{eq:u_h_star}
\end{equation}%

\begin{lemma}\label{lemma:midpoint}
    Let $\Gamma_p$ be a polygon, a closed curve consisting of line segments $\Tilde{\Gamma}_q$, embedded in $\Omega \subset \mathbb{R}^2$ that is composed of 
a union of line segments $\Delta \Gamma_j$, that is ${\Gamma_p = \bigcup_{j=1}^m\Delta \Gamma_j}$, such that each $\Tilde{\Gamma}_q$ consists of an integer number of line segments $\Delta \Gamma_p$. Further, let line segment $\Delta \Gamma_j$ be bounded by vertices $\mathbf{x}_j$ and $\mathbf{x}_{j+1}$, such that 
$\mathbf{x}_{m+1} = \mathbf{x}_{1}$, and let $\mathbf{x}_{j+\frac12}$ be the midpoint of a line segment $\Delta \Gamma_j$, and let $f$ be a $C^2$--smooth function that is defined over an open set 
that includes $\Gamma_p$. Then there exists a $C_1>0$, such that
$$
\bigg| \int_{\Gamma_p} f d \Gamma - \sum_{j = 1}^m f(\mathbf{x}_{j+\frac12}) || \mathbf{x}_{j+1} - \mathbf{x}_j ||  \bigg|\le C_1 h^2.
$$
where $h=\text{max}_{j\in{\{1,...,m}\}}\parallel\mathbf{x}_{j+1} - \mathbf{x}_j\parallel$.

In three dimensions, let $\Gamma_p$ be a closed surface in $\Omega \subset \mathbb{R}^3$, that is composed of a union of surface segments $\Delta \Tilde{\Gamma}_q$. Let $\Gamma_p$ be divided by triangular elements $\Delta \Gamma_j$, such that ${\Gamma_p = \bigcup_{j=1}^m\Delta \Gamma_j}$, let $\mathbf{x}_j$ be the centre of $\Delta \Gamma_j$, then there exists a positive constant $C_2$, such that for each component we have
\[
\left| \int_{\Gamma_p} f({\bf x}) \, d{\bf x} - \sum_{j=1}^{m} \frac{|A_j|}{2} f({\bf x}_j) \right| \leq {C_2}h{^2},
\]
where $h$ is the maximal diameter among all the surface elements over $\Gamma_p$ and $| A_j |/2$ is the area of the (triangular) surface element $\Delta \Gamma_j$.
\end{lemma}

\begin{proposition}
\label{corr2}
    Let $\boldsymbol{u}^{\ast }$ and $\boldsymbol{u}_{h}^{\ast }$, respectively, be given by equation (\ref{eq:u_star}) and (\ref{eq:u_h_star}), and suppose that the numerical quadrature method over the closed surface or contour over $\Gamma_p$ is bounded from above by $C h^r$ for a $C > 0$, $r > 0$. Further let $S$ and $\Tilde{\Gamma}$ be a bounded surface and curve, then there are $C_S > 0$ and $C_{\tilde{\Gamma}}>0$, such that $\parallel \boldsymbol{u}_{h}^*-\boldsymbol{u}^*\parallel_{{\bf L}^2(S)} \leq {C_S}h^{r}$ and $\parallel \boldsymbol{u}_{h}^*-\boldsymbol{u}^*\parallel_{{\bf L}^2(\tilde{\Gamma})} \leq {C_{\tilde{\Gamma}}}h^{r}$, respectively.
\end{proposition} 

 \begin{proof}
     Since ${\bf G}(\boldsymbol{x};\boldsymbol{x}')$ is smooth for any $\boldsymbol{x} \in S$, and hence the requirements of Lemma \ref{lemma:midpoint} are fulfilled, we have $|{\boldsymbol u}_h^* - {\boldsymbol u}^*| \le C_2 h^{r}$ for any $\boldsymbol{x} \in S$. Integration over $S$ gives

\begin{equation*}
\begin{split}  
 \left[\int_{S}| \boldsymbol{u}_{h}^*-\boldsymbol{u}^*|^{2} d{S}\right]^{1/2}\leq \left[\int_{S}C_2^{2}({\boldsymbol{x}})h^{2r} d{S}\right]^{1/2}\le \sup_{\boldsymbol{x} \in S}C_2({\boldsymbol{x}})h^{r} |{S}|^{1/2}=C_S h^{r}.
\end{split}
\end{equation*}
Hence,
\begin{equation}
\parallel \boldsymbol{u}_{h}^*-\boldsymbol{u}^*\parallel _{{\bf L}^{2}(S
)}\leq {C_S}h^{r}.
\label{two}
\end{equation}
and analogously,
\begin{equation}
\parallel \boldsymbol{u}_{h}^*-\boldsymbol{u}^*\parallel _{{\bf L}^{2}(\Tilde{\Gamma}
)}\leq {C_{\tilde{\Gamma}}}h^{r}.
\label{three}
\end{equation}
\end{proof}

\end{document}
